\section{Integral Soul\'e generation at the first irregular prime}
\label{sec:integral}

This section applies McCallum--Sharifi's integral result to Brown's
freeness theorem. It gives a negative answer to integral generation
by the specified Soul\'e generators, already in one of the
characters singled out by $K$-theory.

Let $\mathfrak g=\bigoplus_{m\geq1}\mathfrak g_m$ be the
integral graded Lie algebra of the outer Galois image, for its
weight filtration, in the convention of
\cite[\S9, pp.~24--28]{MS}. Choose the Soul\'e elements
$\sigma_m$ for odd $m\geq3$ with the normalization used there,
and let $\mathfrak h$ be the graded $\Z_p$-Lie subalgebra
generated by their leading classes.

\begin{proposition}\label{prop:integral37}
At $p=37$, the inclusion $\mathfrak h_{32}\subseteq
\mathfrak g_{32}$ is strict. Its cokernel is a nonzero finite
cyclic $37$-group. The character of this graded piece is
$\omega^{32}=\omega^{-4}$.
\end{proposition}
\begin{proof}
McCallum--Sharifi's Theorem~9.7, p.~28, applies at the first
irregular index $m$ when Vandiver holds, the Deligne--Ihara
map is injective through degree $m$, and one of its cup
coefficients $e_{i,m}$ with $i<m/2$ is nonzero.
At $p=37$ the first irregular index is $32$ and Vandiver
holds, as calculated above. Brown's Theorem~1.1 supplies
the Deligne--Ihara injectivity: the rational free Lie map
is injective, so its integral restriction is injective
because a free $\Z_{37}$-Lie algebra is torsionfree.
The required cup coefficient is explicit in
\cite[Corollary~8.3, p.~23]{MS}: the coefficient of
$[x_7,x_{25}]$ is $-1$ modulo $37$, and $7<16$.
Thus the hypotheses of their Theorem~9.7 all hold.
Its conclusion gives
\[
 \mathfrak h_{32}\varsubsetneq\mathfrak g_{32},
 \qquad
 \mathfrak g_{32}=\Z_{37}\bar x_{32}+\mathfrak h_{32}.
\]
The quotient is consequently nonzero and cyclic.
Hain--Matsumoto's rational generation, with Brown's
freeness, gives
$\mathfrak h_{32}\otimes\Q_{37}
 =\mathfrak g_{32}\otimes\Q_{37}$.
These graded pieces are finite $\Z_{37}$-modules in
each weight, so the quotient is finite. It is therefore
$\Z/37^r$ for some $r\geq1$.
The torus acts in weight $32$ by $\omega^{32}$, and
$32\equiv-4\pmod{36}$.
\end{proof}

This application strengthens the formerly conjectural
injectivity input of the published integral obstruction
by using Brown's theorem; the obstruction and its cup
coefficient are those of McCallum and Sharifi.
The argument proves no particular value of $r$.
In particular the filtered relation involving
$B_{32}\bar x_{32}$ has an error in the third term of an
induced filtration. Identifying that error with the
third term of a Lie filtration generated by the chosen
Soul\'e elements would itself require integral generation.
The displayed cyclic quotient is the conclusion justified
by the cited theorem and rational generation.
